\chapter{Precisión numérica y verificación}
\label{cap:precision}

Un método rápido que da el número equivocado no sirve de nada. Este capítulo
trata de cómo saber que el resultado es correcto, que es la parte del trabajo
que más se omite y la que más caro sale omitir.

\section{Las tres fuentes de error}

\begin{enumerate}[leftmargin=1.4em]
\item \textbf{Discretización.} La grilla y la cadena de Markov aproximan un
modelo continuo. Este error existe antes de aplicar ningún método de solución y
es compartido con cualquier alternativa.
\item \textbf{Diferenciación.} Si se usa diferenciación numérica en lugar de
automática, hay un error de orden $h$ (un lado) o $h^2$ (dos lados).
\item \textbf{Truncamiento.} Se resuelve un sistema de horizonte $T$ en lugar de
infinito.
\end{enumerate}

Solo la segunda y la tercera son propias del método. Vale la pena tratarlas por
separado, porque los remedios son distintos.

\section{Verificar contra el método directo}

La prueba más importante, y la que hay que correr siempre que se escribe un
bloque nuevo, es comparar el algoritmo fake news contra el método directo. El
método directo es lento pero no depende de la proposición \ref{prop:principal}:
simplemente perturba y resuelve la transición completa. Si ambos coinciden, el
bloque está bien escrito.

\begin{figure}[htbp]
\centering
\includegraphics[width=\textwidth]{fig05_verificacion.pdf}
\caption{Izquierda: tiempo de cómputo de los dos métodos en el bloque de hogares
de Krusell--Smith con 100 puntos de activos. El método directo escala como $T^2$
y el algoritmo fake news esencialmente como $T$. Derecha: discrepancia máxima
relativa entre ambos. El error se mantiene en el orden de $5\times10^{-9}$, que
es el límite de la diferenciación numérica de dos lados, no un error del
algoritmo.}
\label{fig:verificacion}
\end{figure}

La figura \ref{fig:verificacion} muestra el resultado para el modelo de este
manual. Con diferenciación de dos lados y $h=10^{-4}$, la discrepancia relativa
máxima sobre los cuatro jacobianos ($\J^{A,r}$, $\J^{A,w}$, $\J^{C,r}$,
$\J^{C,w}$) es de $5.5\times 10^{-9}$. El artículo original, usando
diferenciación automática, obtiene coincidencia a precisión de máquina; es decir,
la proposición \ref{prop:principal} es exacta y toda la discrepancia residual
viene de las derivadas numéricas.

\begin{receta}[Verificación mínima de un bloque nuevo]
\begin{enumerate}[leftmargin=1.4em,itemsep=2pt]
\item Resolver el estado estacionario y chequear que las restricciones
contables se cumplan \emph{exactamente} (el presupuesto agregado, la suma de la
distribución igual a uno).
\item Calcular los jacobianos con fake news y con el método directo a un
horizonte corto ($T=30$ o $T=40$, donde el directo todavía es barato) y comparar.
Exigir error relativo menor que $10^{-6}$.
\item Verificar las identidades analíticas que el modelo implique.
\item Recién entonces subir a $T=300$ y resolver el equilibrio general.
\end{enumerate}
\end{receta}

El paso 2 es barato: a $T=40$ el método directo tarda menos de un segundo. No
hay excusa para saltárselo, y saltárselo es la razón número uno por la que la
gente publica resultados equivocados con este método.

\section{Elegir el horizonte de truncamiento}

El horizonte $T$ tiene que ser lo bastante largo para que \emph{tanto el choque
como la respuesta endógena} hayan vuelto esencialmente al estado estacionario.
Son dos condiciones distintas y la segunda es la que se olvida.

\begin{figure}[htbp]
\centering
\includegraphics[width=0.52\textwidth]{fig06_truncamiento.pdf}
\caption{Error de la respuesta del capital en los primeros 60 trimestres,
medido contra una solución de referencia con $T=600$, en función del horizonte
de truncamiento. Con un choque de persistencia $\rho=0.9$ bastan unos 150
trimestres para alcanzar precisión de doce cifras; con $\rho=0.99$ hace falta
bastante más.}
\label{fig:truncamiento}
\end{figure}

La figura \ref{fig:truncamiento} cuantifica el punto. Con $\rho=0.9$ el choque
mismo ya es despreciable hacia el trimestre 100, y el error cae muy rápido. Con
$\rho=0.99$ el choque conserva el 5\% de su magnitud inicial en el trimestre
300, y el error cae mucho más despacio. Modelos con alta persistencia
\emph{interna} --- como los HANK de dos activos, donde la cartera ilíquida se
ajusta lentamente --- son más exigentes todavía.

\begin{receta}[Elegir $T$]
\begin{enumerate}[leftmargin=1.4em,itemsep=1pt]
\item Verificar que el choque mismo sea numéricamente cero en $T$.
\item Verificar que la respuesta endógena de interés sea numéricamente cero en
$T$.
\item Resolver con $T$ y con $1.5T$; si la respuesta en los primeros períodos
cambia de forma apreciable, $T$ es corto.
\end{enumerate}
\end{receta}

\section{Elegir el paso de diferenciación}

Con diferenciación numérica hay un compromiso clásico: $h$ demasiado grande
introduce error de truncamiento de Taylor; $h$ demasiado pequeño amplifica el
error de redondeo. Para diferenciación de dos lados, el óptimo está alrededor de
$h \approx \epsilon^{1/3}\approx 6\times10^{-6}$ en aritmética de doble
precisión, pero en la práctica conviene algo mayor, del orden de $10^{-4}$,
porque el ``ruido'' de la función no es el épsilon de máquina sino la tolerancia
con que convergió el estado estacionario.

\begin{trampa}
Si uno baja $h$ y el jacobiano \emph{empeora} en lugar de mejorar, el
diagnóstico casi seguro es que el estado estacionario está mal convergido: el
ruido del punto base domina la señal de la derivada. La respuesta no es ajustar
$h$ sino converger mejor el estado estacionario.
\end{trampa}

\section{Quiebres, esquinas y lo que la linealización no ve}
\label{sec:quiebres}

Hay una condición que se verifica poco y que invalida el método cuando falla: el
modelo tiene que ser \emph{diferenciable en el estado estacionario}. Las
funciones $\max$, $\min$ y $(\cdot)^+$ son omnipresentes en macrofinanzas y no
son diferenciables en su quiebre.

Hay dos situaciones y conviene no confundirlas.

\paragraph{El estado estacionario está estrictamente dentro de una rama} Si
$\min(1, P/Q)$ se evalúa en un estado estacionario donde $P/Q = 0.36$, la función
es localmente idéntica a $P/Q$ y es suave. No hay problema: basta verificarlo.

\paragraph{El estado estacionario está exactamente en el quiebre} Entonces la
derivada no existe y el método no es aplicable. Hay que recalibrar para salir del
quiebre, o cambiar de método.

Pero hay un tercer caso, más interesante, que es la razón por la que estos
modelos funcionan mucho mejor de lo que uno temería.

\begin{idea}
\textbf{La integración suaviza los quiebres.} Si $x$ es una variable aleatoria
con densidad continua, entonces $\Ex[(x+\mu)^+]$ es una función \emph{derivable
con continuidad} de $\mu$, aunque $(\cdot)^+$ no lo sea. En efecto,
\begin{equation}
\frac{\partial}{\partial\mu}\,\Ex\big[(x+\mu)^+\big] = \Pr(x+\mu > 0),
\end{equation}
que es continua si la densidad lo es. El quiebre individual no sobrevive a la
agregación sobre un continuo.
\end{idea}

\begin{figure}[htbp]
\centering
\includegraphics[width=\textwidth]{fig08_quiebre_esperanza.pdf}
\caption{El faltante de liquidez esperado por unidad de depósitos, con un choque
uniforme en $[-a,a]$. La línea punteada es el caso individual $(-\mu)^+$, con su
quiebre en cero. Las curvas sólidas son la esperanza, que es suave; su derivada
(panel derecho) es la probabilidad de faltante, continua y decreciente. Cuanto
mayor la dispersión $a$, más suave la función.}
\label{fig:quiebre}
\end{figure}

La figura \ref{fig:quiebre} ilustra el punto con la función que aparece en la
parte V. Dos lecturas adicionales que conviene extraer:

\begin{itemize}[leftmargin=1.4em]
\item La función \emph{límite} cuando $a\to 0$ es el caso individual con su
quiebre. Es decir, la suavidad se compra con dispersión: un modelo con muy poca
heterogeneidad idiosincrática está numéricamente cerca de un quiebre, aunque
formalmente sea suave.
\item La agregación suaviza \emph{dos veces}: sobre el choque idiosincrático y
sobre la distribución de agentes. Un agregado sobre bancos con distintos
umbrales es suave aunque cada banco individual tenga un umbral neto.
\end{itemize}

\begin{tutesis}
Esto responde una objeción que seguramente te van a hacer al pasar a dinámico:
``tu modelo está lleno de $\max\{x,0\}$, ¿cómo lo linealizas?''. La respuesta
corta: lo que entra en el modelo no es $\max\{x,0\}$ sino
$\Ex[\max\{x,0\}]$, y esa función es $C^1$ siempre que $F_k$ sea continua ---
que es exactamente el supuesto que tu Anexo IV ya hace. La respuesta larga tiene
una advertencia: la decisión de \emph{colchón voluntario} sí puede estar en una
esquina ($u^c = 0$), y ahí la derivada es cero por un lado y distinta de cero
por el otro. Hay que verificar en qué lado está cada tipo de banco en el estado
estacionario, y reportarlo. La sección \ref{sec:regimenes} lo hace para el
prototipo.
\end{tutesis}
